% GENERATED FILE. Edit canonical lessons, then run npm run generate:books.
% canonical-source-hash: fnv1a64:90be21bcc7f4764e
% canonical-lessons: GE-C16-ich-war, GE-C16-du-warst, GE-C16-er-war, GE-C16-wir-waren, GE-C16-drei-wurzeln, GE-C16-suppletion, GE-C16-haeufig-unregelmaessig, GE-C16-war-practice

\chapter{Three Verbs Under One Roof}
\label{ch:ge-sein-drei-wurzeln}
\hlchaptermodality{19}

\begin{chapteropening}
\textbf{By the end of this chapter:} \emph{I can put the German verb for “to be” into the past, and say why its forms look like three unrelated words — because they are.}

The past of sein is the tidy half of the verb: one stem war-, ordinary endings, and ich war identical to er war. The chapter teaches those four cells one at a time and then answers the question the present tense raised — bin, ist and war come from three different Proto-Indo-European verbs, which is called suppletion, and it survives because the words you never have to guess at are the words analogy never reaches.
\end{chapteropening}

\section[ich war]{ich war — \textquotedblleft{}I was\textquotedblright{}}

\label{lesson:GE-C16-ich-war}



\begin{quote}
You know the simple past is mostly a written tense in German, and that a handful of very common verbs keep it in speech. This is the most common of them.
\end{quote}

\subsection*{You'll want to know: ich war}
\begin{quote}
\textbf{ich war} — \textquotedblleft{}I was.\textquotedblright{}
\end{quote}

Nobody says \emph{ich bin gewesen} in ordinary speech when \emph{ich war} will do. This is the past-tense form Germans actually use, in every region, out loud.

And hear the shape of it. \emph{Bin} has a \emph{b}. \emph{Ist} has an \emph{s} and a \emph{t}. \emph{War} has a \emph{w} and an \emph{r}, and no sound in common with either. That is now three completely unalike stems inside one verb.

\begin{sounds}
\begin{itemize}
\raggedright
  \item \texttt{w-as-v} — German \textbf{w} is English \textbf{v}, so \emph{war} begins like \emph{var}.
  \item \texttt{vowel-a-long} — the \textbf{a} is long and open, the \emph{a} of English \emph{father}.
  \item \texttt{r-final} — the final \textbf{r} softens towards a vowel rather than curling.
\end{itemize}

The result sounds close to English \emph{var}, and nothing like English \emph{war}.
\end{sounds}

\subsection*{Your turn}
Say these aloud:

\begin{itemize}
\raggedright
  \item \textquotedblleft{}ich war\textquotedblright{} — \emph{var}, long open \emph{a}
  \item \textquotedblleft{}ich bin … ich war\textquotedblright{} — now and then
  \item \textquotedblleft{}Ich war müde.\textquotedblright{}
\end{itemize}

\emph{Twice through:} \textquotedblleft{}ich bin, ich war.\textquotedblright{}

\begin{tcolorbox}[breakable,colback=teal!4,colframe=teal!35!black,title={Before you move on}]
Say \textquotedblleft{}I was.\textquotedblright{} (\textbf{Ich war}.) How does the \emph{w} sound? (Like English \textbf{v}.) Is this form used in speech or only in writing? (\textbf{In speech, everywhere} — it is one of the survivors.) How many unlike stems does this verb now have? (\textbf{Three}.) Next: the same form for \emph{du}.
\end{tcolorbox}
\section[du warst]{du warst — \textquotedblleft{}you were\textquotedblright{}}

\label{lesson:GE-C16-du-warst}



\begin{quote}
You have \emph{ich war}. The familiar \textquotedblleft{}you\textquotedblright{} wants the ending it always wants.
\end{quote}

\subsection*{You'll want to know: du warst}
\begin{quote}
\textbf{du warst} — \textquotedblleft{}you were.\textquotedblright{}
\end{quote}

\emph{War} plus \textbf{-st}, with nothing else moved. The same ending that made \emph{bist} out of \emph{bin}, and \emph{gehst} out of \emph{gehen}.

So the past tense of this verb is better behaved than its present: you take one stem, \emph{war-}, and add ordinary endings to it. Everything difficult about \emph{sein} happens in the present.

\begin{sounds}
\begin{itemize}
\raggedright
  \item \texttt{u-long} — \emph{du} keeps its long rounded \textbf{oo}.
  \item \texttt{final-t} — and \emph{warst} ends on the fully released \textbf{t} German always gives you, after the \emph{s}. Say the \emph{s} and the \emph{t} both: \emph{varst}.
\end{itemize}
\end{sounds}

\subsection*{Your turn}
Say these aloud:

\begin{itemize}
\raggedright
  \item \textquotedblleft{}du warst\textquotedblright{} — \emph{varst}, both consonants released
  \item \textquotedblleft{}ich war … du warst\textquotedblright{}
  \item \textquotedblleft{}du bist … du warst\textquotedblright{} — now and then
\end{itemize}

\emph{Twice through:} \textquotedblleft{}ich war, du warst.\textquotedblright{}

\begin{tcolorbox}[breakable,colback=teal!4,colframe=teal!35!black,title={Before you move on}]
Say \textquotedblleft{}you were,\textquotedblright{} to a friend. (\textbf{Du warst}.) Which ending is that? (\textbf{-st}.) Which stem is it added to? (\emph{\textbf{War-}}.) Is the past of this verb regular or irregular in its endings? (\textbf{Regular} — one stem, ordinary endings.) Next: the form for everyone else, and a surprise.
\end{tcolorbox}
\section[er war]{er war — \textquotedblleft{}he was\textquotedblright{}}

\label{lesson:GE-C16-er-war}



\begin{quote}
In the present, \emph{ich} and \emph{er} took different words — \emph{bin} against \emph{ist}. In the past they take the same one.
\end{quote}

\subsection*{You'll want to know: er war}
\begin{quote}
\textbf{er war} — \textquotedblleft{}he was.\textquotedblright{}
\end{quote}

Not a new form. It is \emph{war}, unchanged, the very word you learned for \emph{ich}. And as with \emph{ist}, it covers \textbf{sie war} (\textquotedblleft{}she was\textquotedblright{}), \textbf{es war} (\textquotedblleft{}it was\textquotedblright{}) and \textbf{das war} (\textquotedblleft{}that was\textquotedblright{}) as well.

This is a habit of the German simple past, not a quirk of this verb: the I-form and the he-form are always identical there. \emph{Ich sagte} and \emph{er sagte}. \emph{Ich hatte} and \emph{er hatte}. The present tense never does this; the past always does.

\begin{sounds}
\begin{itemize}
\raggedright
  \item \texttt{r-final} — two softening \emph{r} sounds in a row now, one closing \emph{er} and one closing \emph{war}. Neither curls.
  \item \texttt{vowel-a-long} — and the long open \textbf{a} between them carries the phrase.
\end{itemize}
\end{sounds}

\subsection*{Your turn}
Say these aloud:

\begin{itemize}
\raggedright
  \item \textquotedblleft{}er war, sie war, es war\textquotedblright{}
  \item \textquotedblleft{}Das war gut.\textquotedblright{}
  \item \textquotedblleft{}ich war … er war\textquotedblright{} — the same word twice
\end{itemize}

\emph{Twice through:} \textquotedblleft{}Das war gut.\textquotedblright{}

\begin{tcolorbox}[breakable,colback=teal!4,colframe=teal!35!black,title={Before you move on}]
Say \textquotedblleft{}he was.\textquotedblright{} (\textbf{Er war}.) How is it different from \emph{ich war}? (\textbf{It isn't} — same word.) Does the present do that? (\textbf{No} — \emph{bin} against \emph{ist}.) Say \textquotedblleft{}that was good.\textquotedblright{} (\textbf{Das war gut}.) Next: the plural.
\end{tcolorbox}
\section[wir waren]{wir waren — \textquotedblleft{}we were\textquotedblright{}}

\label{lesson:GE-C16-wir-waren}



\begin{quote}
One more form, and the verb is finished — present and past, all of it.
\end{quote}

\subsection*{You'll want to know: wir waren}
\begin{quote}
\textbf{wir waren} — \textquotedblleft{}we were.\textquotedblright{}
\end{quote}

\emph{War} plus the plural ending \textbf{-en}, the one \emph{wir} takes on every verb. And just as \emph{sind} served three subjects at once, \emph{waren} covers \textbf{sie waren} (\textquotedblleft{}they were\textquotedblright{}) and \textbf{Sie waren} (\textquotedblleft{}you were,\textquotedblright{} politely) with no change.

The one plural form left over is \textbf{ihr wart}, for a group of friends. It takes the plural-\emph{you} \textbf{-t}, and it is the rarest form in the whole verb.

\begin{cousinweb}
English does the identical thing and nobody notices:

\noindent
\begin{tabularx}{\linewidth}{@{}>{\raggedright\arraybackslash}X>{\raggedright\arraybackslash}X>{\raggedright\arraybackslash}X@{}}
\toprule
\textbf{} & \textbf{German} & \textbf{English} \\
\midrule
singular & \emph{war} & \textbf{was} \\
plural & \emph{waren} & \textbf{were} \\
\bottomrule
\end{tabularx}

An \textbf{s} in the singular, an \textbf{r} in the plural, in both languages, inherited from the same ancestor. English \emph{was/were} and German \emph{war/waren} are the same pair of words.
\end{cousinweb}

\subsection*{Your turn}
Say these aloud:

\begin{itemize}
\raggedright
  \item \textquotedblleft{}wir waren, sie waren, Sie waren\textquotedblright{}
  \item \textquotedblleft{}ich war, du warst, er war, wir waren\textquotedblright{}
  \item \textquotedblleft{}war — was; waren — were\textquotedblright{}
\end{itemize}

\emph{Twice through:} \textquotedblleft{}ich war, wir waren.\textquotedblright{}

\begin{tcolorbox}[breakable,colback=teal!4,colframe=teal!35!black,title={Before you move on}]
Say \textquotedblleft{}we were.\textquotedblright{} (\textbf{Wir waren}.) Which English pair matches \emph{war/waren}? (\emph{\textbf{Was/were}}.) What is the plural ending on it? (\textbf{-en}.) Say \textquotedblleft{}they were.\textquotedblright{} (\textbf{Sie waren}.) Next: why one verb needed three different stems.
\end{tcolorbox}
\section[bin, ist, war]{bin, ist, war — three verbs under one roof}

\label{lesson:GE-C16-drei-wurzeln}



\begin{quote}
Say \emph{bin}, \emph{ist}, \emph{war} one after another. They do not look like forms of the same word — and they aren't.
\end{quote}

\begin{cousinweb}
Three separate Proto-Indo-European verbs each handed German a slice of this paradigm, and each of them left cousins in languages you can check:

\noindent
\begin{tabularx}{\linewidth}{@{}>{\raggedright\arraybackslash}X>{\raggedright\arraybackslash}X>{\raggedright\arraybackslash}X@{}}
\toprule
\textbf{forms} & \textbf{from} & \textbf{which also gave} \\
\midrule
\textbf{ist, sind, seid, sein} & *\emph{h\textsubscript{1}es-} — \textquotedblleft{}to be\textquotedblright{} & Latin \emph{est, sunt}; French \emph{est, sont} \\
\textbf{bin, bist} & *\emph{b\textsuperscript{h}uH-} — \textquotedblleft{}to grow, become\textquotedblright{} & English \textbf{be}; Latin \emph{fuī} \\
\textbf{war, waren} & *\emph{wes-} — \textquotedblleft{}to dwell, remain\textquotedblright{} & English \textbf{was, were} \\
\bottomrule
\end{tabularx}

Read the middle row again. \emph{*b\textsuperscript{h}uH-} did not mean \textquotedblleft{}to be\textquotedblright{} at all. It meant \textquotedblleft{}to grow\textquotedblright{} — and the reason English says \emph{be} and German says \emph{bin} is that a verb about growing got promoted into a verb about existing, in both languages, before either of them existed.

The bottom row is the same story with a different verb: \emph{*wes-} meant \textquotedblleft{}to dwell, to remain.\textquotedblright{} Staying somewhere became having been.
\end{cousinweb}

\begin{culture}
\emph{Bin} has a cousin further afield than English. Latin \emph{fuī}, \textquotedblleft{}I was,\textquotedblright{} is from the same \textbf{*\emph{b\textsuperscript{h}uH-}} — and \emph{fuī} is the direct source of Spanish \textbf{fui}, \textquotedblleft{}I was / I went.\textquotedblright{}

So German \emph{bin} and Spanish \emph{fui} are the same ancient word. One landed in the present tense, the other in the past, and the two languages now use it to mean opposite ends of time.
\end{culture}

\subsection*{Your turn}
Say these aloud:

\begin{itemize}
\raggedright
  \item the three roots — \textquotedblleft{}ist from *\emph{h\textsubscript{1}es-}, bin from *\emph{b\textsuperscript{h}uH-}, war from *\emph{wes-}\textquotedblright{}
  \item \textquotedblleft{}ist — est — is\textquotedblright{}
  \item \textquotedblleft{}war — was; waren — were\textquotedblright{}
  \item the cousin abroad — \textquotedblleft{}German bin, Spanish fui, one root\textquotedblright{}
\end{itemize}

\begin{tcolorbox}[breakable,colback=teal!4,colframe=teal!35!black,title={Before you move on}]
How many ancient roots is \emph{sein} built from? (\textbf{Three}.) Which one gives \emph{ist} — and which Latin word proves it? (\textbf{*\emph{h\textsubscript{1}es-}}; Latin \emph{est}.) What did the root behind \emph{bin} originally mean? (\textquotedblleft{}\textbf{To grow, to become}.\textquotedblright{}) Which Spanish word shares it? (\emph{\textbf{Fui}}.) Which root gives English \emph{was} and \emph{were}? (\textbf{*\emph{wes-}}, the same one as \emph{war}.) Next: the name for what this verb is doing.
\end{tcolorbox}
\section[Suppletion]{Suppletion — filling a verb in from another verb}

\label{lesson:GE-C16-suppletion}



\begin{quote}
Three roots, one verb. There is a name for that, and once you have it you will find the same thing in English.
\end{quote}

\begin{grammarlens}[title={Grammar Lens: what suppletion is}]
\begin{quote}
When one verb's forms are filled in from a \textbf{different} verb, that is \textbf{suppletion}.
\end{quote}

\emph{Sein} has no forms of its own for \textquotedblleft{}I am\textquotedblright{} or \textquotedblleft{}I was.\textquotedblright{} It borrowed them, long ago, from two other verbs, and then all three grew together into a single paradigm that native speakers use without ever noticing the seam.
\end{grammarlens}

\begin{cousinweb}
English does the identical thing with its own commonest verbs, and you have never thought twice about it:

\noindent
\begin{tabularx}{\linewidth}{@{}>{\raggedright\arraybackslash}X>{\raggedright\arraybackslash}X>{\raggedright\arraybackslash}X@{}}
\toprule
\textbf{present} & \textbf{past} & \textbf{where the past came from} \\
\midrule
\textbf{go} & \textbf{went} & \emph{wend}, a different verb entirely \\
\textbf{am} & \textbf{was} & the same *\emph{wes-} that gave German \emph{war} \\
\textbf{good} & \textbf{better} & a different word again \\
\bottomrule
\end{tabularx}

\emph{Went} belonged to \emph{wend} — \textquotedblleft{}to make one's way,\textquotedblright{} still alive in \emph{wend your way}. English took it, gave it to \emph{go}, and left \emph{wend} with a regular past of its own. That is precisely what German did to build \emph{sein}.

Suppletion is not a German oddity, and it is not rare. It is simply what languages do to the words they use most — and every language you have looked at in this book has done it to its verb for \textquotedblleft{}to be.\textquotedblright{}
\end{cousinweb}

\subsection*{Your turn}
Say these aloud:

\begin{itemize}
\raggedright
  \item the definition — \textquotedblleft{}one verb's forms, filled in from another verb\textquotedblright{}
  \item \textquotedblleft{}go, went\textquotedblright{} — English doing it
  \item \textquotedblleft{}bin, war\textquotedblright{} — German doing it
\end{itemize}

\begin{tcolorbox}[breakable,colback=teal!4,colframe=teal!35!black,title={Before you move on}]
What is suppletion? (\textbf{One verb's forms filled in from another verb}.) Name the English example. (\emph{\textbf{Go / went}}.) Which verb did \emph{went} actually belong to? (\emph{\textbf{Wend}}.) How many verbs are hiding inside \emph{sein}? (\textbf{Three}.) Next: why this always happens to the commonest words.
\end{tcolorbox}
\section[warum sein unregelmäßig ist]{Why \textquotedblleft{}to be\textquotedblright{} is the messiest word in every language}

\label{lesson:GE-C16-haeufig-unregelmaessig}



\begin{quote}
Every language you have ever met has a broken verb for \textquotedblleft{}to be.\textquotedblright{} That is too consistent to be an accident, and it isn't one.
\end{quote}

\begin{grammarlens}[title={Grammar Lens: frequency protects irregularity}]
\begin{quote}
\textbf{The most-used words are the most irregular.}
\end{quote}

Regularity spreads by \textbf{analogy}. You meet an unfamiliar verb, you do not know its forms, so you guess them from the pattern — the \emph{-e}, \emph{-st}, \emph{-t}, \emph{-en} machine you learned early. Guess often enough and the guess becomes the form: the odd old shape is sanded down and the verb joins the regular crowd.

But you never have to guess at \textquotedblleft{}to be.\textquotedblright{} You hear it a hundred times a day, from the day you start speaking. Every one of those hearings reinforces the strange shape faster than analogy can wear it away.

Rare words get regularised. Common words are protected fossils.
\end{grammarlens}

\begin{culture}
This also tells you where the irregularity will sit inside a verb, and \emph{sein} proves it. Its present tense — the part you say constantly — is the wreckage: three roots, five unlike forms. Its past tense, which you say far less often, has quietly become well-behaved: one stem \emph{war-}, ordinary endings, \emph{warst} built exactly like \emph{bist} was not.

So the mess is not spread evenly. It clusters where the traffic is.

The practical consequence is a comfort. There is no long list of verbs this strange, because a verb has to be \emph{this} common to survive being this strange — and there are only a handful of words in any language that common.
\end{culture}

\subsection*{Your turn}
Say these aloud:

\begin{itemize}
\raggedright
  \item the rule — \textquotedblleft{}the most-used words are the most irregular\textquotedblright{}
  \item why — \textquotedblleft{}you never have to guess at them\textquotedblright{}
  \item \textquotedblleft{}bin, bist, ist\textquotedblright{} against \textquotedblleft{}war, warst, waren\textquotedblright{} — mess, then order
\end{itemize}

\begin{tcolorbox}[breakable,colback=teal!4,colframe=teal!35!black,title={Before you move on}]
State the rule. (\textbf{The most-used words are the most irregular}.) What spreads regularity? (\textbf{Analogy} — guessing from the pattern.) Why does analogy never reach \textquotedblleft{}to be\textquotedblright{}? (\textbf{You never have to guess at it}.) Which half of \emph{sein} is the tidy half? (\textbf{The past} — one stem, ordinary endings.) Next: the chapter together.
\end{tcolorbox}
\section[Practice]{Practice — was, were, and why}

\label{lesson:GE-C16-war-practice}



\begin{quote}
No new words. Present and past, side by side, and then the account of how one verb ended up like this.
\end{quote}

\begin{grammarlens}[title={Grammar Lens: the past, assembled}]
Six forms, one stem, four endings, no surprises. Said in order:

\begin{itemize}
\raggedright
  \item \textbf{ich war}, \textbf{du warst}, \textbf{er war} — and the first and third are the same word, which the present tense would never allow;
  \item \textbf{wir waren}, \textbf{ihr wart}, \textbf{sie waren}.
\end{itemize}

Compare that with the present, where five different shapes came from three different verbs. The past of \emph{sein} is the part of the paradigm that behaved.
\end{grammarlens}

\subsection*{The exchange}
\noindent
\begin{tabularx}{\linewidth}{@{}>{\raggedright\arraybackslash}X>{\raggedright\arraybackslash}X@{}}
\toprule
\textbf{German} & \textbf{English} \\
\midrule
\emph{Wo warst du?} & Where were you? \\
\emph{Ich war in Berlin.} & I was in Berlin. \\
\emph{Wie war es?} & How was it? \\
\emph{Es war gut. Aber ich bin müde.} & It was good. But I am tired. \\
\bottomrule
\end{tabularx}

\emph{Wo} (\textquotedblleft{}where\textquotedblright{}) and \emph{in} are read straight off the English here; neither is asked of you yet.

\subsection*{Your turn}
Say these aloud:

\begin{itemize}
\raggedright
  \item \textquotedblleft{}ich war, du warst, er war, wir waren, ihr wart, sie waren\textquotedblright{}
  \item the pairs — \textquotedblleft{}ich bin / ich war\textquotedblright{}, \textquotedblleft{}wir sind / wir waren\textquotedblright{}
  \item the three roots — \textquotedblleft{}ist, bin, war — three verbs\textquotedblright{}
  \item the whole exchange, both voices
\end{itemize}

\emph{Twice through:} \textquotedblleft{}Wie war es? Es war gut.\textquotedblright{}

\begin{tcolorbox}[breakable,colback=teal!4,colframe=teal!35!black,title={Before you move on}]
Run the six past forms. Which two are identical? (\emph{\textbf{Ich war}} and \emph{\textbf{er war}}.) Name the three roots this verb is built from, by their German forms. (\emph{\textbf{Ist}}, \emph{\textbf{bin}}, \emph{\textbf{war}}.) What is the technical name for that? (\textbf{Suppletion}.) Why has nothing tidied this verb up? (\textbf{It is too common — you never guess at it, so analogy never reaches it}.) Next: the past tense that is built on this verb rather than out of it.
\end{tcolorbox}
